% page 627 du fac-simile (image p-626 du scan)
% bloc 167.6 x 250.6 mm ; marge gauche 34.4 mm
\pgc{12.700mm}{0.000mm}{
\l{}
\l{\cel{54}619}
\l{}
\l{\sou{vistar}. (\sou{trans.})\cel{2}Aplikar (sur akto) la formulo qua konstatas ke}
\l{\cel{3}lu esas verifikita da la persono o da la organismo di qua}
\l{\cel{3}la signaturo esas necesa por igar olu valida. - EIS.}
\l{}
\l{\sou{vito}. (\sou{bot.}) Planto kun stipo lignatra,qua produktas la beri}
\l{\cel{3}de qui facesas vino. - L.\cel{1}\sou{vitis vinifera}. - EFIS.}
\l{}
\l{\sou{vit}alismo.\cel{2}(\sou{filoz.)}\cel{2}Doktrino segun qua la fenomeni fiziolo-}
\l{\cel{3}giala havas kauzo specala, ne-sama kam olta di la fenomeni}
\l{\cel{3}psikologiala. - DEFIRS.}
\l{}
\l{\sou{vitelo}. La parto flava di la ovo. - EFIS.}
\l{}
\l{\sou{vitro}. Materio harda, frajila, diafana, quan onu obtenas per}
\l{\cel{3}la kunfuzo di sablo silicoza kun potaso e natroxo. - EFIS.}
\l{}
\l{\sou{vitralio}. Panelo de vitro-plaki, asemblita fako-formacante, e}
\l{\cel{3}(maxim-multa-kaze) dekorita ye pikturi. - F.}
\l{}
\l{\sou{vitriolo}. Nomo vulgara di SO4 H2. - DEFIS.}
\l{}
\l{\sou{viulo}. (\sou{olim)}. Nomo di muzikili diversa ; kun kordi ed arketo,}
\l{\cel{3}forme di granda violino kun kin, sis o sep kordi, ne-egale}
\l{\cel{3}grosa. - DEFIRS.}
\l{}
\l{\sou{vivar}.\cel{2}(\sou{netrans.}) I. (\sou{che la animalo}). Posedar e praktikar la}
\l{\cel{3}funcioni \textquotedbl{}nutro\textquotedbl{}, \textquotedbl{}genito\textquotedbl{}. - II. (\sou{che la homo}). Posedar e}
\l{\cel{3}praktikar la funcioni \textquotedbl{}nutro\textquotedbl{}, \textquotedbl{}genito\textquotedbl{}, \textquotedbl{}raciono\textquotedbl{} ed \textquotedbl{}arbi-}
\l{\cel{3}trio\textquotedbl{}. - EFIS.}
\l{}
\l{\sou{vivaca}. Rapida ed animata pri lua ago-maniero.- EFIS.}
\l{}
\l{\sou{vivipara}.\cel{2}Qua parturas ja vivanta sua yuni (konstraste kun}
\l{\cel{3}\textquotedbl{}ovipara\textquotedbl{}). -\cel{2}EF.}
\l{}
\l{\sou{vizar}. (\sou{trans.}) Regardar tre atencoze (lor sua lanso o jus ante}
\l{\cel{3}sua lanso di objekto ad la punto quan onu volas atingar).-}
\l{}
\l{}
\l{\sou{vizajo}. La facio di la kapo (\sou{che la homo}). - EFI.}
\l{}
\l{\sou{vizelo}. (\sou{zool.}) Varietato ek la genero \textquotedbl{}martro\textquotedbl{}, kun korpo svelta,}
\l{\cel{3}longa, falva, falveso plu klara subventre, muzelo pinto-forma.}
\l{\cel{3}- L.\cel{1}\sou{mustela vulgaris}. - DE.}
\l{}
\l{\sou{viziero}. Parto avana di la kasko, quan onu abasis por shirmar}
\l{\cel{3}la vizajo, e tra qua onu povis vidar e respirar per greto}
\l{\cel{3}streta. - Peco ek ledro muntita an-avan kasqueto por shirmar}
\l{\cel{3}la okuli kontre la suno-radii e kontre la pluvoguti. - DEFIS.}
\l{}
\l{\sou{viziono}. Iluziono per qua onu kredas vidar ulu od ulo qua ne}
\l{\cel{3}esas prezenta koram onua okuli. - DEFIS.}
\l{}
\l{\sou{viziro}. (en\cel{1}\sou{Turkia olima}). Oficiro ek la\cel{2}konselio (konsilantaro)}
\l{\cel{3}di la sultano. - DEFIRS.}
}
